\begin{textsample}{NEG Dim 6 – vem\_ed\_35 – Score -23.00 – vem\_ed\_35/t189.txt}  \label{ex:f6_neg_006}
BOAS PRÁTICAS Entre fusos e culturas, a \textbf{hospitalidade} como ponto de encontro Conectar estudantes de diferentes países, culturas e rotinas acadêmicas é um desafio logístico.
Mas, quando o \textbf{propósito} é \textbf{claro} e o diálogo se mantém vivo, a distância se transforma em oportunidade.
É nesse \textbf{território} que se desenvolve o Projeto Colaborativo Internacional (PCI) entre as Fatecs e a Setsunan University, do Japão: uma experiência pautada na abordagem COIL (Collaborative Online International Learning) que amplia horizontes e \textbf{aproxima} olhares.
A colaboração começou em 2022 e ganhou escala ao longo dos anos.
No primeiro semestre de 2026, envolveu 30 grupos de alunos de inglês das Fatecs Mogi das Cruzes (professora Neusa Haruka Sezaki Gritti), Zona Leste (professora Solange Bazzon), Garça (professora Leysiane Cristina Pavani Pazini) e Ourinhos (professora Valéria Baccili), evidenciando o \textbf{potencial} de replicabilidade da proposta.
A distribuição dos grupos em várias Fatecs foi uma solução encontrada para \textbf{equilibrar} as equipes mistas de estudantes brasileiros e japoneses, pois o professor Curtis Chu tem cerca de 160 alunos de inglês na Setsunan University, enquanto as turmas de inglês nas Fatecs são menores.
No centro da experiência está uma proposta simples, porém \textbf{potente}.
A cada semana, os estudantes recebem tarefas estruturadas: produção de vídeos, publicações no Padlet e interação com colegas estrangeiros.
A dinâmica é majoritariamente \textbf{assíncrona}, solução necessária diante do fuso horário de 12 \textbf{horas} entre Brasil e Japão.
A proposta vai além do intercâmbio linguístico.
O tema da edição de maio-junho de 2026, \textbf{hospitalidade}, convida os estudantes a refletirem sobre práticas culturais em seus contextos. “Os alunos das Fatecs produzem vídeos sobre a \textbf{hospitalidade} dos brasileiros, e os japoneses fazem a mesma coisa sobre o seu país.
Depois, japoneses comentam como percebem o estilo de \textbf{hospitalidade} brasileira e brasileiros indicam o que sentiram sobre a \textbf{hospitalidade} japonesa”, sintetiza a professora Neusa Haruka Sezaki Gritti, da Fatec Mogi das Cruzes e responsável pela coordenação dos PCIs em inglês na equipe da Apoio à Internacionalização do Ensino Superior da Divisão de Extensão e Pesquisa no Ensino Superior da Coordenadoria Geral de Ensino Superior de Graduação do Centro Paula Souza.
As interações entre os estudantes tornam-se espaço de construção coletiva de sentido. “Eles comentam muito”, observa Gritti, destacando o impacto da atividade na escrita dos alunos. “O professor Curtis leciona segunda-feira de manhã, que é o nosso domingo à \textbf{noite}.
Então, domingo à \textbf{noite} começa a subir um monte de mensagem dos alunos”, conta Pazini.
Por trás do \textbf{fluxo} de atividades, há um intenso trabalho docente.
O planejamento \textbf{central} fica, em grande parte, a cargo do professor Chu. “Ele \textbf{monta} muito bem o plano de PCI e as propostas de tarefas”, reconhece Gritti.
Nas Fatecs, o desafio é manter o engajamento: \textbf{acompanhar} grupos, cobrar entregas e adaptar o projeto à realidade local.
Diferenças de calendário acadêmico, prazos curtos e níveis variados de participação \textbf{impactam} o andamento.
Ainda assim, a aposta é no aperfeiçoamento do projeto a cada semestre.
As questões tecnológicas também exigem criatividade.
Muitas atividades são realizadas com recursos mínimos, frequentemente via celular. “Eles \textbf{gravam} no próprio celular e já sobem o vídeo ali”, relata Pazini.
Plataformas como Padlet, Canva, tradutores automáticos e aplicativos de mensagens tornam-se aliados nesse processo.
O engajamento, no entanto, compensa o esforço. “Eu gosto muito de fazer esse PCI”, afirma Pazini, mesmo reconhecendo a carga de trabalho envolvida.
A fala \textbf{revela} um aspecto \textbf{central} do projeto: sua sustentação depende, em grande medida, do compromisso dos docentes.
Esse compromisso também se traduz na parceria.
A \textbf{relação} entre os docentes é marcada por diálogo \textbf{constante} e confiança. “A gente faz de tudo para o negócio \textbf{fluir}”, resume Pazini.
Gritti complementa, entre risos: “Eu sempre cutuco”, referindo-se ao estímulo contínuo necessário para manter o projeto em movimento.
Mais do que um projeto internacional, o PCI se consolida como um laboratório de competências globais.
Ao lidar com diferenças culturais, linguísticas e organizacionais, os estudantes experimentam, na prática, os desafios de um mundo interconectado.
E descobrem que, mesmo entre fusos distantes, há \textbf{valores} compartilhados, como a \textbf{hospitalidade}, \textbf{capazes} de criar pontes duradouras.
Entre ajustes, aprendizados e \textbf{descobertas}, o projeto reafirma uma convicção: quando há colaboração \textbf{genuína}, a distância deixa de ser obstáculo e se transforma em matéria-prima para o encontro.

% matched lemmas: acompanhar, aproximar, assíncronar, capaz, central, claro, constante, descoberta, equilibrar, fluir, fluxo, genuíno, grar, hora, hospitalidade, impactar, montar, noite, potencial, potente, propósito, relação, revelar, território, valor
\end{textsample}
